\documentclass[10pt]{beamer}
\usepackage{booktabs}
\usepackage{pgfplots}
\pgfplotsset{compat=1.18}
\usepackage{xcolor}
\usepackage{array}

\usetheme{Madrid}
\usecolortheme{seahorse}

\title{Understanding the Differences Between MORTAR and GEAR}
\subtitle{Dataset Structure, Schema and Distribution Analysis}
\date{\today}

\begin{document}

%% Slide 1 ─ Title
\begin{frame}
  \titlepage
\end{frame}

%% Slide 2 ─ Purpose
\begin{frame}{Purpose of the Comparison}
  \begin{itemize}
    \item Four prediction approaches are evaluated on the GEAR building:
          rule-based (PBRR, AnyBURL), tree-based (GBT), and graph neural networks (CompGCN / LLM).
    \item All approaches depend critically on the structure and schema of the input graph.
    \item MORTAR (US commercial buildings, Brick schema) was used for development and training.
    \item GEAR (Singapore lab building) is the held-out transfer target.
    \item \textbf{This analysis is model-independent.} Findings apply to all four approaches.
    \item Objective: measure \emph{how} GEAR differs from MORTAR so that schema gaps,
          structural mismatches, and vocabulary shifts can be addressed at the data level.
  \end{itemize}
\end{frame}

%% Slide 3 ─ Dataset Overview
\begin{frame}{Dataset Overview}
  \begin{table}
    \centering
    \begin{tabular}{lrr}
      \toprule
      \textbf{Property} & \textbf{MORTAR} & \textbf{GEAR} \\
      \midrule
Buildings & 44 & 1 \\
Total entities & 19,493 & 1,762 \\
Norm.\ triples & 59,428 & 9,309 \\
Unique predicates & 8 & 12 \\
Unique RDF classes & 85 & 46 \\
Target edges & 21284 & 656 \\
Non-target context edges & 38144 & 8653 \\
REC$\to$Brick remaps & 0 & 0 \\
Inverse folds applied & 828 & 623 \\
      \bottomrule
    \end{tabular}
  \end{table}
\end{frame}

%% Slide 4 ─ Target-Relation Distribution
\begin{frame}{Target-Relation Distribution}
  \begin{table}
    \centering
    \begin{tabular}{lrrrr}
      \toprule
      \textbf{Relation} & \multicolumn{2}{c}{\textbf{MORTAR}} & \multicolumn{2}{c}{\textbf{GEAR}} \\
      \cmidrule(lr){2-3}\cmidrule(lr){4-5}
                        & Count & \% & Count & \% \\
      \midrule
\texttt{brick:feeds} & 4450 & 20.9\% & 74 & 11.3\% \\
\texttt{brick:hasPart} & 4790 & 22.5\% & 387 & 59.0\% \\
\texttt{brick:hasPoint} & 12044 & 56.6\% & 195 & 29.7\% \\
      \midrule
\textbf{Total} & 21284 & 100\% & 656 & 100\% \\
      \bottomrule
    \end{tabular}
  \end{table}
  \vspace{0.5em}
  \begin{itemize}
    \small
    \item MORTAR is dominated by \texttt{hasPoint}; GEAR has a much larger \texttt{hasPart} share.
    \item GEAR \texttt{hasPart} includes IFC collection membership (\texttt{rec:includes}),
          absent in MORTAR.
  \end{itemize}
\end{frame}

%% Slide 5 ─ Predicate and Schema Differences
\begin{frame}{Predicate and Schema Differences}
  \begin{columns}[t]
    \column{0.48\textwidth}
    \textbf{Common in MORTAR only} (top)\par
    \vspace{0.3em}
    \begin{itemize}
      \item \texttt{owl:imports} (44)
      \item \texttt{brick:area} (43)
    \end{itemize}
    \column{0.48\textwidth}
    \textbf{Common in GEAR only} (top)\par
    \vspace{0.3em}
    \begin{itemize}
      \item \texttt{\textless{}https://w3id.org/rec\textgreater{}:locatedIn} (601)
      \item \texttt{\textless{}https://w3id.org/rec\textgreater{}:hasPart} (558)
      \item \texttt{\textless{}http://qudt.org/schema/qudt\textgreater{}:hasUnit} (114)
      \item \texttt{\textless{}https://w3id.org/rec\textgreater{}:levelNumber} (9)
    \end{itemize}
  \end{columns}
  \vspace{0.5em}
  \begin{itemize}
    \item Predicate vocabulary Jaccard similarity: 0.4286
    \item Predicate distribution JSD: 0.2161 (0 = identical)
    \item GEAR uses REC namespace predicates; 0 remapped to Brick equivalents.
    \item \texttt{brick:hasLocation} appears 0 times in GEAR, 0 in MORTAR.
    \item Inverse-relation usage: 623 folds in GEAR vs 828 in MORTAR.
  \end{itemize}
\end{frame}

%% Slide 6 ─ Entity and Class Differences
\begin{frame}{Entity and Class Differences}
  \begin{columns}[t]
    \column{0.48\textwidth}
    \textbf{Top MORTAR classes}\par
    \begin{itemize}
      \item \texttt{brick:HVAC\_Zone} (2137)
      \item \texttt{brick:VAV} (2136)
      \item \texttt{brick:Zone\_Air\_Temperature\_Sensor} (2085)
      \item \texttt{brick:Room} (2025)
      \item \texttt{brick:Zone\_Air\_Temperature\_Setpoint} (1631)
    \end{itemize}
    \column{0.48\textwidth}
    \textbf{Top GEAR classes}\par
    \begin{itemize}
      \item \texttt{brick:Equipment} (587)
      \item \texttt{\textless{}https://w3id.org/rec\textgreater{}:Room} (548)
      \item \texttt{brick:HVAC\_System} (330)
      \item \texttt{brick:Status} (35)
      \item \texttt{brick:Emergency\_Alarm} (33)
    \end{itemize}
  \end{columns}
  \vspace{0.4em}
  \begin{itemize}
    \item Class vocabulary Jaccard: 0.0738 --- most GEAR classes absent from MORTAR.
    \item GEAR-only classes include spatial types (Room, Floor, Building) absent in MORTAR training.
    \item Untyped entities: MORTAR 106 / GEAR 46
    \item After REC$\to$Brick mapping, spatial GEAR classes gain Brick ancestor embeddings.
  \end{itemize}
\end{frame}

%% Slide 7 ─ Ontology Category Distribution
\begin{frame}{Ontology Category Distribution}
  \begin{table}
    \centering
    \small
    \begin{tabular}{lrrrr}
      \toprule
      \textbf{Category} & \textbf{MORTAR} & \textbf{M\%} & \textbf{GEAR} & \textbf{G\%} \\
      \midrule
Equipment & 3008 & 15.5 & 607 & 35.4 \\
Point & 11248 & 58.0 & 218 & 12.7 \\
Location & 4275 & 22.1 & 0 & 0.0 \\
Building & 0 & 0.0 & 0 & 0.0 \\
System & 0 & 0.0 & 0 & 0.0 \\
Collection & 0 & 0.0 & 332 & 19.3 \\
Other & 0 & 0.0 & 0 & 0.0 \\
Unrecognised & 856 & 4.4 & 559 & 32.6 \\
      \bottomrule
    \end{tabular}
  \end{table}
\end{frame}

%% Slide 8 ─ Relation Type-Pattern Differences
\begin{frame}{Relation Type-Pattern Differences}
  \textbf{\texttt{brick:hasPart}} subject$\to$object category pairs\par
  \begin{table}
    \small
    \begin{tabular}{llrr}
      \toprule
      Subject & Object & MORTAR & GEAR \\
      \midrule
Location & Location & 4050 & 0 \\
Equipment & Equipment & 641 & 8 \\
Collection & Equipment & 0 & 379 \\
Equipment & Other & 99 & 0 \\
      \bottomrule
    \end{tabular}
  \end{table}
  \begin{itemize}
    \item MORTAR \texttt{hasPart}: Equipment$\to$Equipment (HVAC containment).
    \item GEAR \texttt{hasPart}: Location$\to$Location (IFC/BIM spatial hierarchy) +
          Collection$\to$* (grouping).
    \item These are semantically distinct uses of the same predicate.
  \end{itemize}
\end{frame}

%% Slide 9 ─ Naming and URI Differences
\begin{frame}{Naming and URI Differences}
  \begin{table}
    \centering
    \begin{tabular}{lrr}
      \toprule
      \textbf{Metric} & \textbf{MORTAR} & \textbf{GEAR} \\
      \midrule
Avg local-name length & 62.6 & 52.4 \\
Avg tokens per name & 9.02 & 8.75 \\
\% with underscores & 75.8\% & 94.7\% \\
\% with hyphens & 5.1\% & 33.1\% \\
\% with numeric IDs & 99.5\% & 97.6\% \\
      \bottomrule
    \end{tabular}
  \end{table}
  \begin{itemize}
    \item MORTAR entities use \texttt{bldgN:} namespace (anonymised building IDs).
    \item GEAR entities use a full IFC filename as namespace prefix.
    \item Name-similarity features (if used) benefit MORTAR's consistent conventions.
    \item GEAR contains \texttt{rdfs:label} annotations; MORTAR entities rely on URI tokens only.
  \end{itemize}
\end{frame}

%% Slide 10 ─ Graph Structure and Context
\begin{frame}{Graph Structure and Available Context}
  \begin{table}
    \centering
    \begin{tabular}{lrr}
      \toprule
      \textbf{Metric} & \textbf{MORTAR} & \textbf{GEAR} \\
      \midrule
Avg entity degree & 5.155 & 7.321 \\
Max entity degree & 2137 & 587 \\
Connected components & 5 & 3 \\
Largest component & 19384 & 1754 \\
Remaining triples (target removed) & 38144 & 8653 \\
Isolated after removal & 20 & 0 \\
Target-subj retaining context (\%) & 99.6\% & 100.0\% \\
      \bottomrule
    \end{tabular}
  \end{table}
  \begin{itemize}
    \item Non-target context in GEAR is almost entirely \texttt{rdf:type} and \texttt{brick:hasLocation}.
    \item Removing target edges leaves many GEAR entities with only their type triple.
  \end{itemize}
\end{frame}

%% Slide 11 ─ Main Findings
\begin{frame}{Main Findings}
  \textbf{Confirmed Differences}
  \begin{itemize}
    \item \texttt{hasPart} semantics differ fundamentally: HVAC containment (MORTAR) vs.\
          IFC spatial hierarchy + collection membership (GEAR).
    \item GEAR class vocabulary is largely disjoint from MORTAR (Jaccard 0.0738).
    \item GEAR uses REC predicates (0 remapped); MORTAR uses Brick natively.
    \item GEAR non-target context is sparse: mostly \texttt{rdf:type} + \texttt{hasLocation} only.
    \item \texttt{hasPart} share: 59.0\% in GEAR vs.\
          22.5\% in MORTAR.
  \end{itemize}
  \vspace{0.3em}
  \textbf{Likely Effects on Graph-Prediction Methods}
  \begin{itemize}
    \item Models trained on MORTAR \texttt{hasPart} patterns will not generalise to GEAR's
          IFC/collection use.
    \item Sparse context reduces embedding quality for GNN-based approaches.
    \item Schema heterogeneity affects any method that uses predicate identity as a feature.
  \end{itemize}
  \vspace{0.3em}
  \textbf{Issues Requiring Further Testing}
  \begin{itemize}
    \item Whether REC$\to$Brick remapping fully resolves predicate vocabulary mismatch.
    \item Whether additional Brick ancestor expansion compensates for unseen class distributions.
  \end{itemize}
\end{frame}

%% Slide 12 ─ Recommended Data-Level Next Steps
\begin{frame}{Recommended Data-Level Next Steps}
  \begin{itemize}
    \item \textbf{Schema harmonisation:} ensure all GEAR triples are mapped to Brick before
          evaluation; validate completeness of REC$\to$Brick predicate and class mapping.
    \item \textbf{Correct malformed URIs:} check for double-prefix Brick URIs and
          colon-in-local-name patterns in both datasets.
    \item \textbf{Separate evaluation by relation:} report \texttt{feeds}, \texttt{hasPart},
          and \texttt{hasPoint} metrics independently to isolate domain-shift effects.
    \item \textbf{Separate evaluation by entity-type group:} e.g., report results
          for Equipment-only, Location-only, and Collection subsets of GEAR.
    \item \textbf{Expand training data:} include buildings with IFC/BIM spatial containment
          hierarchies and REC-annotated data if available.
    \item \textbf{Document GEAR custom classes:} entities lacking a Brick equivalent
          should be explicitly labelled as out-of-vocabulary rather than treated as unknown.
    \item \textbf{Create a GEAR validation subset:} a small labelled subset stratified by
          relation type enables calibrated threshold selection without using the test set.
  \end{itemize}
\end{frame}

\end{document}
